\chapter{Results and Discussion}
\label{chap:results}

\chapterabstract{This chapter reports the main results, ablations, qualitative examples, and limitations.}

\section{Main Results}
\placeholder{Present the primary quantitative results first. Explain the trend, not just the numbers.}

\begin{table}[t]
  \centering
  \caption{Main result table placeholder.}
  \label{tab:main-results}
  \begin{tabular}{@{}lccc@{}}
    \toprule
    Method & Metric 1 & Metric 2 & Metric 3 \\
    \midrule
    Baseline & -- & -- & -- \\
    Proposed & -- & -- & -- \\
    \bottomrule
  \end{tabular}
\end{table}

\section{Ablation Studies}
\placeholder{Isolate the effect of each contribution. Keep ablations tied to the claims made in the introduction.}

\section{Qualitative Results}
\placeholder{Show representative successes and failures. Qualitative figures should make the method's behavior inspectable.}

\section{Limitations}
\placeholder{Discuss failure cases, data bias, compute cost, assumptions, and evaluation gaps. Strong limitations make the thesis more credible.}

\section{Discussion}
\placeholder{Connect the results back to the motivation. Explain what the experiments imply for the broader research problem.}
